% !TEX root = ../../main.tex
\subsection{安全目标}

结合外卖订单通信场景和上述威胁模型，FoodShield 的安全目标主要包括匿名性、认证性、机密性、完整性、可追责性、可审计性和抗滥用能力。系统通过不同安全机制分别支撑上述目标，使外卖订单通信过程不仅能够完成业务沟通，而且能够满足基本的信息安全要求。

\begin{table}[H]
\centering
\caption{FoodShield 系统安全目标与对应机制}
\renewcommand{\arraystretch}{1.7}
\setlength{\tabcolsep}{8pt}
\zihao{-4}

\begin{tabular}{p{3.0cm}p{5.0cm}p{5.0cm}}
\toprule
\textbf{安全目标} & \textbf{防护对象} & \textbf{对应机制} \\
\midrule

匿名性
&
用户真实身份、用户 PID、跨订单关联信息
&
PID 内部化；骑手端不可见 PID；前端字段最小化展示
\\

认证性
&
订单通信入口、订单参与资格
&
HMAC-SM3 订单 Token；绑定 orderID、PID、timestamp、nonce
\\

机密性
&
订单聊天内容、数据库消息记录
&
WebSocket/WSS 传输保护；SM4 消息加密存储；管理员端默认不看明文
\\

完整性
&
历史通信记录、消息顺序、消息数量
&
SM3 消息摘要；Merkle Tree；Merkle Root 快照比对
\\

可追责性
&
投诉纠纷、恶意骚扰、异常订单
&
用户端与骑手端均可按 orderID 发起申请；管理员条件溯源
\\

可审计性
&
管理员查询、验证、溯源等高权限操作
&
管理员认证；audit\_logs 记录操作时间、目标订单和操作结果
\\

抗滥用能力
&
机器注册、垃圾订单、匿名身份滥用
&
用户名与口令认证；机器识别；注册频率限制
\\

隐私最小化
&
用户备注、标签、通信明文、身份映射关系
&
管理员端仅展示订单摘要、消息哈希、Merkle Root 和验证状态
\\

\bottomrule
\end{tabular}
\end{table}




其中，匿名性和可追责性是本系统设计中的一组核心平衡目标。完全匿名虽然能够增强隐私保护，但可能导致恶意行为难以治理；完全实名虽然便于追责，但会增加用户身份暴露风险。因此，FoodShield 采用可控匿名设计：日常通信中隐藏真实身份，异常场景下通过管理员权限校验和审计记录完成条件溯源。

认证性是系统防止越权通信的关键。用户进入订单通信通道前，需要提交订单号和订单 Token，服务器根据订单绑定关系重新计算 Token 并判断其合法性。由于 Token 由订单号、PID、时间戳和随机数共同参与生成，攻击者在不知道订单认证密钥的情况下难以伪造合法凭证。

完整性是系统可信审计能力的基础。FoodShield 不仅保存通信记录，还对每条消息生成摘要，并将消息摘要组织为 Merkle Tree。系统通过保存 Merkle Root 快照，使后续审计过程中能够验证历史通信记录是否发生变化。如果任意消息内容、消息顺序或消息数量被修改，重新计算的 Root 将与历史快照不一致，从而暴露篡改行为。

可审计性主要用于约束管理员权限。管理员虽然可以处理异常订单和执行溯源，但其操作必须被记录。系统通过审计日志保存管理员身份、操作时间、目标订单、操作类型和操作结果，使管理员行为具备可追踪性，从而降低审计能力本身带来的隐私风险。
